% !TeX program = xelatex
% ============================================================
% pz/pz.tex — Пояснительная записка (ГОСТ 19.404-79)
%
% Главный содержательный документ курсового проекта: описывает
% назначение, постановку задачи, архитектуру и функционирование
% программы, обоснование выбора средств и ожидаемые технико-
% экономические показатели. По составу заменяет ВКР в курсовом
% комплекте (та же роль «пояснительного» документа, но по ЕСПД
% ГОСТ 19.404-79, а не по ГОСТ 7.32).
%
% Жёлтые блоки \hseFill{…} — места, которые нужно заменить своим
% текстом. Данные (ФИО, название, коды) приходят из настроек
% проекта через макросы \hse… (common/meta.tex).
% ============================================================
\documentclass[12pt,a4paper]{extarticle}
\input{../common/preamble}
\newcommand{\docname}{Explanatory Note}
% ЕСПД-код пояснительной записки — тип «81» (ГОСТ 19.103-77).
\newcommand{\doccode}{\hseDocCodeBase\ 81 01-1}
\newcommand{\doccodelu}{\hseDocCodeBase\ 81 01-1-ЛУ}
\input{../common/commands}
\input{../common/styles}
\begin{document}
\sloppy
\input{../common/title_template}


% ============================================================
%  ENGLISH BODY (project.lang == en)
%  §2.3.4.2: an applied English project may present its ЕСПД documentation
%  in English as a translation of GOST 19. The structure of GOST 19.404-79
%  is preserved section-by-section; references follow IEEE style.
%  Fill the yellow \hseFill placeholders with your own text.
% ============================================================

% ============================================
% ANNOTATION
% ============================================
\section*{ANNOTATION}
\addcontentsline{toc}{section}{ANNOTATION}

The explanatory note is a document containing information about the algorithms used, the operating principles and other information related to the operation of the program.

This explanatory note \enquote{\hseProjectName} contains the following sections: \enquote{Introduction}, \enquote{Purpose and Field of Application}, \enquote{Technical Characteristics}, \enquote{Expected Technical and Economic Indicators}, \enquote{Sources Used During Development}.

The \enquote{Introduction} section states the name of the program and the conventional designation of the development topic, as well as the basis for the development.

The \enquote{Purpose and Field of Application} section presents the purpose of the program and a brief characterisation of its field of application.

The \enquote{Technical Characteristics} section contains the following subsections: statement of the task for the development of the program; description of the algorithm and functioning of the program; description and justification of the choice of the method for organising input and output data; description and justification of the choice of technical and software resources.

The \enquote{Expected Technical and Economic Indicators} section states the anticipated demand and the economic advantages of the development compared with domestic and foreign models or analogues.

\clearpage

% ============================================
% LIST OF STANDARDS USED
% ============================================
\noindent
The software document \enquote{Explanatory Note} is formatted in accordance with the requirements of the ESPD (Unified System of Program Documentation):

\begin{enumerate}[label=\arabic*), leftmargin=1.25cm]
\item GOST 19.101-77 Types of programs and program documents\cite{gost19101}.
\item GOST 19.103-77 Designations of programs and program documents\cite{gost19103}.
\item GOST 19.104-78 Basic inscriptions\cite{gost19104}.
\item GOST 19.105-78 General requirements for program documents\cite{gost19105}.
\item GOST 19.106-78 Requirements for program documents produced by printing\cite{gost19106}.
\item GOST 19.404-79 Explanatory note. Requirements for content and formatting\cite{gost19404}.
\end{enumerate}

Changes to this explanatory note are made in accordance with GOST\,19.603-78\cite{gost19603}, GOST\,19.604-78\cite{gost19604}.

\clearpage

% ============================================
% CONTENTS
% ============================================
\hypersetup{linkcolor=black}
\tableofcontents
\hypersetup{linkcolor=blue}
\thispagestyle{fancy}
\clearpage

% ============================================
% 1. INTRODUCTION
% ============================================
\section{INTRODUCTION}

\subsection{Name of the program}

% The names are substituted automatically: the full name comes from the
% project title, the English and short names come from the project settings.
The name of the program is \enquote{\projectname}.

The name of the program in English is \enquote{\hseProjectNameEn}.

\subsection{Conventional designation of the development topic}

The conventional designation of the development topic is \enquote{\hseProjectNameShort}.

\subsection{Basis for the development}

% Hint: refer to the curriculum of the study programme and to the topic of the
% course project approved by the academic supervisor.
The basis for the development is the bachelor's curriculum of the study field \hseFill{code and name of the study field} and the topic of the course project approved by the academic supervisor of the \enquote{\hseSpecialization} educational programme.

\clearpage

% ============================================
% 2. PURPOSE AND FIELD OF APPLICATION
% ============================================
\section{PURPOSE AND FIELD OF APPLICATION}

\subsection{Purpose of the program}

\subsubsection{Functional purpose}

% Hint: list the main functions the program performs.
The functional purpose of the program is \hseFill{main functional purpose of the program}. This work implements \hseFill{list the key components and functions}.

\subsubsection{Operational purpose}

% Hint: describe by whom and under what conditions the program is operated.
This program can be used for \hseFill{operation scenarios}. The target audience of the project is \hseFill{list the target users}.

\subsection{Brief characterisation of the field of application}

% Hint: 1--2 paragraphs about the purpose and the subject area in which the
% program is applied and which tasks it solves.
The program is intended for use in \hseFill{subject area and systems where the program is applied}. It can be used both as \hseFill{scenario 1} and as \hseFill{scenario 2}.

\clearpage

% ============================================
% 3. TECHNICAL CHARACTERISTICS
% ============================================
\section{TECHNICAL CHARACTERISTICS}

\subsection{Statement of the task for the development of the program}

% Hint: formulate the task being solved and give the functional requirements
% (you may refer to the requirements from the Statement of Work, prefix ТЗ-Ф-01...).
This work is developed for \hseFill{purpose of developing the program}. The functional requirements for the program are given in the Statement of Work and detailed below.

\begin{hseExample}
Example list of tasks to be solved:
\begin{itemize}[nosep, leftmargin=1.25cm]
    \item \hseFill{task 1};
    \item \hseFill{task 2};
    \item \hseFill{complete the list of tasks}.
\end{itemize}
\end{hseExample}

\subsection{Description of the algorithm and functioning of the program}

% Hint: this is the central section of the explanatory note. Describe the
% architecture, the key algorithms, data schemas and interaction between
% components. Diagrams (component, sequence, class) and pseudocode are welcome.
\subsubsection{Program architecture}

\hseFill{Describe the architecture: the components, their purpose and interaction. Provide an architecture diagram.}

\subsubsection{Description of functioning}

\hseFill{Describe the main operating scenarios of the program: how requests are processed, what the key algorithms and data flows are.}

\subsection{Description and justification of the choice of the method for organising input and output data}

% Hint: formats, protocols, encodings, data schemas; why they were chosen.
\hseFill{Describe the formats and protocols of the input and output data, the data storage schema, and justify the choice.}

\subsection{Description and justification of the choice of technical and software resources}

% Hint: languages, frameworks, DBMS, infrastructure; justification of the choice.
\hseFill{List the languages, frameworks, libraries, DBMS and infrastructure used; justify the choice of each resource.}

\clearpage

% ============================================
% 4. EXPECTED TECHNICAL AND ECONOMIC INDICATORS
% ============================================
\section{EXPECTED TECHNICAL AND ECONOMIC INDICATORS}

\subsection{Anticipated demand}

% Hint: describe the potential users and the anticipated demand.
The developed program may be of interest to \hseFill{target audience}. The potential users are \hseFill{list the potential users}.

\subsection{Economic advantages of the development compared with domestic and foreign models or analogues}

% Hint: briefly describe the advantages of the development over analogues
% (a comparison table may be used).
\hseFill{Describe the economic and functional advantages of the development compared with existing analogues.}

\clearpage
% The source pages are printed without the bottom stamp (as in real
% document sets) --- only the sheet number and the document code at the top.
\pagestyle{appendix}

% ============================================
% SOURCES USED DURING DEVELOPMENT
% ============================================
\section*{SOURCES USED DURING DEVELOPMENT}
\addcontentsline{toc}{section}{SOURCES USED DURING DEVELOPMENT}

% Hint: add the sources on the technologies of your project BEFORE the GOSTs,
% in IEEE style, for example:
% \bibitem{fastapi}
% FastAPI. Accessed: Mar.~14, 2026. [Online]. Available: \url{https://fastapi.tiangolo.com}
\renewcommand{\refname}{}
\begin{thebibliography}{99}
\bibitem{gost19101}
\emph{GOST 19.101-77. Types of Programs and Program Documents}. Unified System of Program Documentation. Moscow, Russia: IPK Izdatelstvo Standartov, 2001.

\bibitem{gost19103}
\emph{GOST 19.103-77. Designations of Programs and Program Documents}. Unified System of Program Documentation. Moscow, Russia: IPK Izdatelstvo Standartov, 2001.

\bibitem{gost19104}
\emph{GOST 19.104-78. Basic Inscriptions}. Unified System of Program Documentation. Moscow, Russia: IPK Izdatelstvo Standartov, 2001.

\bibitem{gost19105}
\emph{GOST 19.105-78. General Requirements for Program Documents}. Unified System of Program Documentation. Moscow, Russia: IPK Izdatelstvo Standartov, 2001.

\bibitem{gost19106}
\emph{GOST 19.106-78. Requirements for Program Documents Produced by Printing}. Unified System of Program Documentation. Moscow, Russia: IPK Izdatelstvo Standartov, 2001.

\bibitem{gost19404}
\emph{GOST 19.404-79. Explanatory Note. Requirements for Content and Formatting}. Unified System of Program Documentation. Moscow, Russia: IPK Izdatelstvo Standartov, 2001.

\bibitem{gost19603}
\emph{GOST 19.603-78. General Rules for Making Changes}. Unified System of Program Documentation. Moscow, Russia: IPK Izdatelstvo Standartov, 2001.

\bibitem{gost19604}
\emph{GOST 19.604-78. Rules for Making Changes to Program Documents Produced by Printing}. Unified System of Program Documentation. Moscow, Russia: IPK Izdatelstvo Standartov, 2001.
\end{thebibliography}

\clearpage

% ============================================
% LIST OF TERMS AND DEFINITIONS
% ============================================
\section*{LIST OF TERMS AND DEFINITIONS}
\addcontentsline{toc}{section}{LIST OF TERMS AND DEFINITIONS}

% Hint: add the terms of the subject area and the technologies of your project.
% Order --- alphabetical.
\begin{enumerate}[label=\arabic*., leftmargin=1.25cm]
\item \textbf{\hseFill{term}}~--- \hseFill{definition of the term}.
\item \textbf{\hseFill{another term}}~--- \hseFill{definition of the term}.
\end{enumerate}

\clearpage


% ============================================
% ЛИСТ РЕГИСТРАЦИИ ИЗМЕНЕНИЙ
% ============================================
\input{../common/change_log}

\end{document}
